%! The Fundamental Theorem of Linear Algebra — Unit 3, Lesson 3.6 — Exit Ticket (Answer Key)
\documentclass[10pt]{article}
\usepackage{linalg-article}
\usepackage{linalg-key}   % pulls in -boxes; do NOT also load -boxes
\newcommand{\TallMath}[1]{$\displaystyle #1\rule[-1.4em]{0pt}{3.2em}$}
\newcommand{\xx}{\mathbf{x}}
\newcommand{\yy}{\mathbf{y}}
\newcommand{\svec}{\mathbf{s}}
\newcommand{\zero}{\mathbf{0}}
\newcommand{\T}{\mathsf{T}}

\begin{document}

\pageheader{Unit 3, Lesson 3.6}{Exit Ticket --- Answer Key}
\namedateperiod

\vspace{0.15in}

No notes. Show your reasoning.

\begin{enumerate}[leftmargin=1.8em, itemsep=15pt]
  \item \textbf{The big picture.} A matrix $A$ is $4\times 3$ with rank $r=2$. Give the four dimensions and
        the room ($\mathbb{R}^m$ or $\mathbb{R}^n$) each subspace lives in, then name the two perpendicular
        pairs (Part~2).
        \[
          \dim C(A)=\ans{$2$}\ \ \dim C(A^{\T})=\ans{$2$}\ \
          \dim N(A)=\ans{$1$}\ \ \dim N(A^{\T})=\ans{$2$}
        \]
        \ansline{$m=4$, $n=3$. $C(A)$ and $N(A^{\T})$ live in $\mathbb{R}^m=\mathbb{R}^4$; $C(A^{\T})$ and $N(A)$ live in $\mathbb{R}^n=\mathbb{R}^3$. Perpendicular pairs: in $\mathbb{R}^3$, $C(A^{\T})\perp N(A)$; in $\mathbb{R}^4$, $C(A)\perp N(A^{\T})$. (Checks: $2+1=3=n$, $2+2=4=m$.)}
  \item \textbf{Check a right angle.} For $A=\begin{bmatrix} 2 & 1 \\ 4 & 2 \end{bmatrix}$, the nullspace
        vector is $\svec=\begin{bsmallmatrix} -1\\2 \end{bsmallmatrix}$ and the top row is $(2,1)$. Compute
        $(2,1)\cdot\svec$ and say what it shows about the nullspace and the row space.
        \ansline{$(2,1)\cdot\svec=(2)(-1)+(1)(2)=-2+2=0$. A dot product of $0$ means perpendicular, so the nullspace vector $\svec$ is perpendicular to the row $(2,1)$ --- i.e.\ $N(A)\perp C(A^{\T})$ (the row space is spanned by that row).}
  \item \textbf{Why?} In one or two sentences, use ``$A\xx$ is computed row by row'' to explain why
        $A\xx=\zero$ forces $\xx$ to be perpendicular to every row of $A$.
        \ansline{Each entry of $A\xx$ is a row of $A$ dotted with $\xx$. So $A\xx=\zero$ makes every one of those dot products $0$, and a dot product of $0$ means the vectors are perpendicular --- hence $\xx$ is perpendicular to every row.}
\end{enumerate}

\begin{teachernote}
Sort responses into ``got it / paired across rooms / checked only one vector.'' Item~1 is the Part~1
discriminator (four dimensions \emph{and} the right room --- note $\dim N(A^{\T})=m-r=2$ uses the row count
$m=4$). Item~2 is the concrete Part~2 check: one dot product, and the reading that dot $=0$ means the two
subspaces are perpendicular. Item~3 is the conceptual crux --- the row-by-row reason. If many pair subspaces
across rooms or state the wrong dimension for $N(A^{\T})$, re-draw the two-room big picture before the Unit~3
test.
\end{teachernote}

\end{document}
